\section{Results}

% ============================================================== R1
\subsection{Data collection and study design}
We designed a workflow that joins data collection, model collection and four kinds of evaluation (Fig.~\ref{fig:design}a). Seven test sets were prepared from six datasets, \rn{n.images} images in total (Fig.~\ref{fig:design}b and Table~\ref{tab:datasets}). IDRiD and DDR are the training sources of the main tables, because both have official splits and almost every published number refers to one of them. TJDR, split into its standard-field and 133$^\circ$ ultra-widefield parts, FGADR, Retinal-Lesions and DiaRetDB1 are independent sources: no model in the main tables was trained on them, and every checkpoint is applied to their test sets as it is. The sources differ in how lesions were drawn. Following the analysis of lesion count and lesion area by Playout and Cheriet\cite{playout2024crossdatasetgeneralization} we assigned each source to one of four granularity tiers, fine, mixed, coarse and region consensus, and we average results only within a tier.

A fixed search of three bibliographic databases identified 1,389 records; 764 remained after a Boolean gate and de-duplication, 112 entered the full-text queue, and screening of the priority list added 24 methods and 2 datasets to a snowballed corpus, for a ledger of 86 records (Fig.~\ref{fig:design}d). From it we registered \rn{n.models} model variants in five families and recorded for each a reproducibility level from R0 (no code) to R4 (training code, weights, splits and environment) (Fig.~\ref{fig:design}c and Table~\ref{tab:models}). \rn{n.tier1} variants form tier~1 and enter the main comparisons: general segmentation networks (U-Net\cite{ronneberger2015unet}, DeepLabv3\cite{chen2017deeplabv3,chen2018deeplabv3plus}, HRNet\cite{wang2020hrnet}, Swin-Unet\cite{cao2022swinunet}, FCT\cite{tragakis2023fct}, H2Former\cite{he2023h2former}), the self-configuring nnU-Net\cite{isensee2021nnunet}, the fundus-specific M2MRF\cite{liu2023m2mrf} and HACDR-Net\cite{xu2024hacdrnet}, and eight foundation models (SAM\cite{kirillov2023sam}, SAM~2\cite{ravi2024sam2}, SAM~3\cite{carion2025sam3}, MedSAM\cite{ma2024medsam}, SAMed\cite{zhang2023samed}, SAM-Adapter\cite{chen2023samadapter}, SAM2-UNet\cite{xiong2024sam2unet} and RETFound\cite{zhou2023retfound}). \rn{n.tier2} variants in tier~2 are added as time allows and \rn{n.tier3} in tier~3 are audited from their papers only.

Each model can appear under four evaluations that are never mixed in one column. The paper audit records the numbers as published and where they came from. Author weights are released checkpoints scored on official and independent test sets. Author recipes are the authors' pipelines rerun from scratch. Unified retraining applies one protocol, called \Bstd{} below, to the tier-1 models with complete training code. Two test scenarios are reported throughout, as in the template study\cite{lu2026tcrepitopenaturemethods}: the internal test, which is the official test set of the training source, and the independent test. The primary metric is the exact AUPR of each lesion class over all field-of-view pixels of a test set, averaged over the four classes (\mAUPR); Dice, intersection over union and lesion-level recall by lesion size are secondary. Finally, six factors are varied one at a time on one to five representative models (Fig.~\ref{fig:design}e).

\begin{figure}[p]
\centering
\widefig{fig1_design}
\caption{\textbf{Data collection and study design.} \textbf{a}, Workflow: the task, the seven test sources, the registered model variants, the four kinds of evaluation, and the two test scenarios with their metrics. \textbf{b}, Images per source and frozen split (train\,/\,validation\,/\,test). Colour gives the annotation granularity tier; shade separates the three parts of a split. \textbf{c}, Registered model variants by family and reproducibility level; the number of tier-1 variants is written after each bar. \textbf{d}, Literature search and screening. \textbf{e}, Factors varied one at a time, with the experiment code used in Methods.}
\label{fig:design}
\end{figure}

\begin{table*}[t]
\centering\tablefont
\caption{\textbf{Datasets, frozen splits and roles.} All eight sources carry pixel- or region-level labels for microaneurysms (MA), hemorrhages (HE), hard exudates (EX) and soft exudates (SE). Granularity decides which test sets may be averaged: averages are taken only within one tier. Split is train\,/\,validation\,/\,test in images; official test sets are never changed. ``Maps'' marks sources whose probability maps we can release.}
\label{tab:datasets}
\begin{adjustbox}{max width=\textwidth}
\begin{threeparttable}
\begin{tabular}{@{}l l l l l l c r c@{}}
\toprule
\colh{Source} & \colh{Ref.} & \colh{Images} & \colh{Resolution} & \colh{FOV} & \colh{Annotation} & \colh{Granularity} & \colh{Split} & \colh{Maps} \\
\midrule
\grouprow{9}{Training sources (internal test = official test set)}
IDRiD & \cite{porwal2018idrid} & 81 \inlinebar{0.14} & 4288$\times$2848 & 50$^\circ$ & pixel contours, expert review & \tierbadge{fine}{Fine} & 44 / 10 / 27 & \yes \\
DDR & \cite{li2019ddr} & 757 \inlinebar{1.32} & variable & 45$^\circ$ & pixel masks & \tierbadge{fine}{Fine} & 383 / 149 / 225 & \yes \\
\grouprow{9}{Independent test sources (never used to train a main-table model)}
TJDR, standard field & \cite{mao2023tjdr} & 257 \inlinebar{0.45} & 2048$\times$2048 & 35--50$^\circ$ & polygons, 2 graders + adjudicator & \tierbadge{fine}{Fine} & 172 / 30 / 55 & \yes \\
TJDR, ultra-widefield & \cite{mao2023tjdr} & 304 \inlinebar{0.53} & 3912$\times$3912 & 133$^\circ$ & polygons, 2 graders + adjudicator & \tierbadge{fine}{Fine} & 209 / 37 / 58 & \yes \\
MAPLES-DR & \cite{maplesdr2024scientificdata} & 198 \inlinebar{0.34} & 1440$\times$960–2304$\times$1536 & 45$^\circ$ & pixel masks, 7 retinologists; MA/HE/EX AI-assisted & \tierbadge{fine}{Fine} & \na\ / 0 / 60 & \yes \\
FGADR & \cite{zhou2021fgadr} & 1,842 \inlinebar{3.20} & 1280$\times$1280 & n.s. & pixel masks, 2 residents + senior & \tierbadge{mixed}{Mixed} & 1,096 / 193 / 553 & \no \\
Retinal-Lesions & \cite{wei2020retinallesionslesionnet} & 1,593 \inlinebar{2.77} & 896$\times$896 & n.s. & ellipses and polygons, $\geq$3 of 45 graders & \tierbadge{coarse}{Coarse} & 951 / 164 / 478 & \no \\
DiaRetDB1 & \cite{kauppi2007diaretdb1} & 89 \inlinebar{0.15} & 1500$\times$1152 & 50$^\circ$ & regions with confidence, 4 experts & \tierbadge{region}{Region} & \na\ / 28 / 61 & \tbd \\
\bottomrule
\end{tabular}
\begin{tablenotes}[flushleft]\scriptsize
\item FOV, field of view; n.s., not stated in the source paper. IDRiD: 10 of the 54 official training images are held out for validation in unified retraining; author-recipe retraining uses all 54. DiaRetDB1: ground truth is the mean confidence of four experts thresholded at 0.75; the 28 official training images are used only to choose thresholds. Retinal-Lesions: split by patient; grey-level 127 regions are ignored. MAPLES-DR: labels drawn at $1500\\times1500$ on the field-of-view crop of MESSIDOR images and placed back at the original resolution; the 138 training images are not used, the official 60 test images are; uncertain red (MA or HE) and bright (EX, SE or drusen) lesions are ignored for the two classes they could belong to. The FGADR licence forbids redistribution of derived data; maps for Retinal-Lesions are not released either.
\end{tablenotes}
\end{threeparttable}
\end{adjustbox}
\end{table*}

\begin{table*}[t]
\centering\tablefont
\caption{\textbf{Models entering the retraining tracks.} 38 variants are registered: 17 in tier~1 (main tables), 8 in tier~2 (if time allows) and 13 in tier~3 (paper audit only; Supplementary Table~2). Reproducibility level: R0, no code; R1, model definition only; R2, inference code and weights; R3, training code; R4, training code, weights, splits and environment. Tracks: A, author weights; B, unified retraining of the official code; C, re-implementation from the paper, reported in a separate column.}
\label{tab:models}
\begin{adjustbox}{max width=\textwidth}
\begin{threeparttable}
\begin{tabular}{@{}l c l c c c l l@{}}
\toprule
 & & & \multicolumn{3}{c}{\colh{Track}} & & \\
\cmidrule(lr){4-6}
\colh{Model} & \colh{Tier} & \colh{Reproducibility} & \colh{A} & \colh{B} & \colh{C} & \colh{Code status} & \colh{Use in this study} \\
\midrule
\grouprow{8}{General segmentation and self-configuring baselines}
\famchip{general}\; U-Net & 1 & \rlevel{3} {\scriptsize R3} & \no & \yes & \no & reference implementation & B-std \\
\famchip{general}\; DeepLabv3 & 1 & \rlevel{3} {\scriptsize R3} & \no & \yes & \no & torchvision, ResNet-101 & B-std \\
\famchip{general}\; HRNet-W48 & 1 & \rlevel{3} {\scriptsize R3} & \no & \yes & \no & official, W48 & B-std \\
\famchip{auto}\; nnU-Net & 1 & \rlevel{4} {\scriptsize R4} & \no & \yes & \no & official, 2D; default and ResEnc-M & own entry point \\
\famchip{general}\; Swin-Unet & 1 & \rlevel{3} {\scriptsize R3} & \no & \yes & \no & official & B-std, 224 patches \\
\famchip{general}\; FCT & 1 & \rlevel{3} {\scriptsize R3} & \no & \yes & \no & official PyTorch port & B-std, 384 patches \\
\famchip{general}\; H2Former & 1 & \rlevel{3} {\scriptsize R3} & \no & \yes & \no & official & deferred: exceeds 40\,GB at 960 \\
\famchip{general}\; UNet++ & 2 & \rlevel{3} {\scriptsize R3} & \no & \yes & \no & official & tier 2 \\
\famchip{general}\; UPerNet & 2 & \rlevel{3} {\scriptsize R3} & \no & \yes & \no & official & tier 2 \\
\grouprow{8}{Fundus-specific, multi-scale CNN}
\famchip{fundus}\; M2MRF & 1 & \rlevel{4} {\scriptsize R4} & \yes & \yes & \no & official; weights released & A, E1r, B-std \\
\famchip{fundus}\; HACDR-Net & 1 & \rlevel{3} {\scriptsize R3} & \no & \yes & \no & official; no weights & E1r, B-std \\
\famchip{fundus}\; WSRFNet & 2 & \rlevel{1} {\scriptsize R1} & \yes & \no & \yes & model and test code only & A if weights arrive; C \\
\famchip{fundus}\; L-Seg & 2 & \rlevel{1} {\scriptsize R1} & \no & \no & \yes & Caffe prototxt and loss layer & C \\
\famchip{fundus}\; SAA & 2 & \rlevel{0} {\scriptsize R0} & \no & \no & \yes & repository empty & C \\
\famchip{fundus}\; PMCNet & 2 & \rlevel{3} {\scriptsize R3} & \no & \yes & \no & official, TF\,1.13; no weights & tier 2 \\
\grouprow{8}{Fundus-specific, relation or transformer}
\famchip{fundus}\; RTNet & 2 & \rlevel{0} {\scriptsize R0} & \no & \no & \yes & no code; image + vessel input & C \\
\grouprow{8}{Foundation models}
\famchip{foundation}\; SAM & 1 & \rlevel{2} {\scriptsize R2} & \no & \yes & \no & official inference & zero-shot, oracle prompt \\
\famchip{foundation}\; SAM 2 & 1 & \rlevel{3} {\scriptsize R3} & \no & \yes & \no & official & zero-shot, oracle prompt \\
\famchip{foundation}\; SAM 3 & 1 & \rlevel{3} {\scriptsize R3} & \no & \yes & \no & official; gated weights & text prompt \\
\famchip{foundation}\; MedSAM & 1 & \rlevel{3} {\scriptsize R3} & \no & \yes & \no & official & oracle box, full fine-tuning \\
\famchip{foundation}\; SAMed & 1 & \rlevel{3} {\scriptsize R3} & \no & \yes & \no & official & PEFT (LoRA) \\
\famchip{foundation}\; SAM-Adapter & 1 & \rlevel{3} {\scriptsize R3} & \no & \yes & \no & official & PEFT \\
\famchip{foundation}\; SAM2-UNet & 1 & \rlevel{3} {\scriptsize R3} & \no & \yes & \no & official & PEFT \\
\famchip{foundation}\; RETFound + UPerNet & 1 & \rlevel{3} {\scriptsize R3} & \no & \yes & \no & official encoder; UPerNet head added & encoder + decoder \\
\famchip{foundation}\; PerSAM & 2 & \rlevel{2} {\scriptsize R2} & \no & \yes & \no & official inference & one-shot, tier 2 \\
\bottomrule
\end{tabular}
\begin{tablenotes}[flushleft]\scriptsize
\item \famchip{general} general segmentation \quad \famchip{auto} self-configuring \quad \famchip{fundus} fundus-specific \quad \famchip{foundation} foundation model. B-std, unified retraining protocol (Methods). PEFT, parameter-efficient fine-tuning. Status as of the code audit of 2--4 October 2026; author replies are due by 31 October 2026.
\end{tablenotes}
\end{threeparttable}
\end{adjustbox}
\end{table*}


% ============================================================== R2
\subsection{The numbers in the literature}
We first asked what stands behind the numbers that papers report. Of the \rn{funnel.dr} fundus-specific methods in the registry (17 dedicated architectures and 3 adaptations of foundation models to this task), \rn{funnel.code} have any author code, \rn{funnel.train} have training code and \rn{funnel.full} (M2MRF) comes with weights, splits and an environment description (Fig.~\ref{fig:literature}a). Authors of methods without code or weights were contacted with a fixed template; methods that could not be run stay in the denominator.

For every method paper in the ledger (\rn{lit.n.methods} method papers and one prior cross-dataset benchmark\cite{playout2024crossdatasetgeneralization}) we extracted the experimental setup and the comparison tables, \rn{lit.base.n} baseline AUPR values on IDRiD and DDR in all (Supplementary Table~3). The official split is the rule: \rn{lit.idrid.official} of the \rn{lit.idrid.n} papers that use IDRiD and \rn{lit.ddr.official} of the \rn{lit.ddr.n} that use DDR state it, and the others give no split or merge the official test images into training. What the papers leave out is how the metric was computed. Of \rn{lit.impl.n} papers whose main metric is AUPR, \rn{lit.impl.notstated} give no implementation, \rn{lit.impl.fixed} integrate over a fixed grid of 11 or 33 thresholds, which is the IDRiD challenge protocol carried on by the evaluation code released with L-Seg and M2MRF, \rn{lit.impl.exact} use the exact average precision and \rn{lit.impl.perimage} averages over images; none states that the evaluation is restricted to the field of view, and \rn{lit.pooled.stated} say whether pixels are pooled over the test set or scores are averaged over images (Fig.~\ref{fig:literature}b). Baseline numbers travel between papers. Of the baseline values, \rn{lit.base.rerun}\% were produced by the paper that reports them, \rn{lit.base.copied1}\% are identical to the value in the method's own paper, \rn{lit.base.copied2}\% are identical to a value that an earlier paper had already copied or produced, and \rn{lit.base.unverified}\% cite a source outside the ledger or none. Only \rn{lit.rerun.all} papers state that every baseline was rerun, and \rn{lit.rerun.notstated} say nothing about where their comparison numbers come from. As a consequence one method carries a range of values: L-Seg is reported on IDRiD by \rn{lit.lseg.n} papers with \mAUPR{} from \rn{lit.lseg.lo} to \rn{lit.lseg.hi}, DeepLabv3+ by \rn{lit.deeplab.n} papers from \rn{lit.deeplab.lo} to \rn{lit.deeplab.hi} and U-Net by \rn{lit.unet.n} papers from \rn{lit.unet.lo} to \rn{lit.unet.hi}. Copied values agree to the last digit; the spread comes from the papers that reran a method under their own settings (Fig.~\ref{fig:literature}c). The audit also surfaced copied tables whose lesion columns are transposed with respect to the source, a baseline table reproduced verbatim from an earlier paper of the same group without attribution, and a Dice value carried into an AUPR column (Supplementary Table~3).

The metric implementation alone moves the numbers. The evaluation code released with M2MRF, which later work reuses, computes AUPR from 11 thresholds; at threshold 0 every pixel with a non-zero probability is counted as positive, so the recall is 1 and the precision equals the prevalence, which lifts the area under the curve. On the released IDRiD checkpoint this function returns \rn{e1.m2mrf.IDRiD.fn}; after the probabilities are stored as 16-bit maps the same function returns \rn{e1.m2mrf.IDRiD.stored}, and the exact pooled AUPR inside the field of view is \rn{e1.m2mrf.IDRiD.exact}. On DDR the three values are \rn{e1.m2mrf.DDR.fn}, \rn{e1.m2mrf.DDR.stored} and \rn{e1.m2mrf.DDR.exact} (Fig.~\ref{fig:literature}d).
\begin{simtext}
Applied to the probability maps of all retrained models, the 11-threshold implementation raised \mAUPR{} on IDRiD by \rn{impl.mean} on average and by up to \rn{impl.max}, by an amount that depends on how each model distributes its low probabilities. \rn{impl.moved} of \rn{n.main} models changed rank (Kendall $\tau$ = \rn{impl.tau}; Fig.~\ref{fig:literature}e).
\end{simtext}

\begin{figure}[t]
\centering
\widefig{fig2_literature}
\caption{\textbf{The numbers in the literature.} \textbf{a}, From registered fundus-specific methods to a reproduced number; the first four stages follow the reproducibility levels of the registry. \textbf{b}, Share of method papers by split and by AUPR implementation (fixed thresholds, a grid of 11 or 33 operating points), and provenance of the baseline AUPR values on IDRiD and DDR in their comparison tables: produced by the reporting paper, identical to the method's own paper, identical to a value already carried by an earlier paper, or citing a source outside the ledger. \textbf{c}, \mAUPR{} on IDRiD reported for one method in different papers (dots; bar, range; $n$, number of papers; methods reported by at least three papers) and the value after unified retraining (diamond). \textbf{d}, The same probability maps scored with the exact AUPR and with the 11-threshold implementation; the two shaded rows are the released M2MRF checkpoints, for which the authors' function on the float output is also shown. \textbf{e}, Rank of the retrained models under the two implementations.}
\label{fig:literature}
\end{figure}

% ============================================================== R3
\subsection{Performance of the original models}
Released checkpoints were scored with three numbers each: the value printed in the paper, the authors' evaluation function, and our protocol (Fig.~\ref{fig:original}a and Table~\ref{tab:original}). For M2MRF the authors' function reproduced the repository values to the last printed digit on both datasets (\rn{e1.m2mrf.IDRiD.fn} and \rn{e1.m2mrf.DDR.fn}), which shows that moving the code to a current framework did not change its output; the paper reports the mean of three runs (\rn{e1.m2mrf.IDRiD.paper} and \rn{e1.m2mrf.DDR.paper}). Under our protocol the same checkpoints reach \rn{e1.m2mrf.IDRiD.exact} and \rn{e1.m2mrf.DDR.exact}, a gap of \rn{e1.m2mrf.IDRiD.gap} and \rn{e1.m2mrf.DDR.gap} that is almost entirely due to the threshold-0 term of the 11-threshold curve; restricting the evaluation to the field of view and keeping overlapping labels each change the third decimal only (Fig.~\ref{fig:original}b).

Rerunning the authors' recipe is the direct test of whether released code reproduces a reported number. A first rerun of M2MRF on IDRiD, with an effective batch of 1 instead of the authors' 4, reached \rn{e1r.first.fn} with the authors' metric and \rn{e1r.first.exact} with ours, well below the released checkpoint.
\begin{simtext}
With the authors' batch size, M2MRF came within \rn{e1r.M2MRF.IDRiD.delta} of the reported value on IDRiD over three seeds and within \rn{e1r.M2MRF.DDR.delta} on DDR in one run, and HACDR-Net, for which no weights are released, within \rn{e1r.HACDR-Net.IDRiD.delta} and \rn{e1r.HACDR-Net.DDR.delta} in one run per source (Fig.~\ref{fig:original}c). On the independent test sets the author weights lost most of their internal performance, with the same ordering of sources as for the retrained models below (Fig.~\ref{fig:original}d).
\end{simtext}

\begin{figure}[t]
\centering
\widefig{fig3_original}
\caption{\textbf{Performance of the original models.} \textbf{a}, \mAUPR{} of each released checkpoint on the official test set of its training source: reported in the paper, authors' function on the checkpoint, and this study's protocol. The number on the right is the difference between the last two. \textbf{b}, Decomposition of that difference for the two M2MRF checkpoints. \textbf{c}, Retraining with the authors' recipe: reported value (line) and individual seeds (dots); the cross is the first rerun at an effective batch of 1. \textbf{d}, Per-class AUPR of the author weights on all seven test sets under this study's protocol. Outlined cells are internal tests.}
\label{fig:original}
\end{figure}

\begin{table*}[t]
\centering\tablefont
\caption{\textbf{Original models: three numbers per released checkpoint, and retraining with the authors' recipe.} \textbf{a}, Each checkpoint is scored on the official test set of its training source with the number printed in the paper, the authors' evaluation function (11 thresholds, whole image, overwritten labels) on the float output, the same function on the stored 16-bit maps, and this study's protocol (exact pooled AUPR inside the field of view, overlapping labels). \textbf{b}, The authors' pipeline rerun from scratch; mean $\pm$ s.d. over seeds, last iteration, no checkpoint selection.}
\label{tab:original}
\begin{adjustbox}{max width=\textwidth}
\begin{threeparttable}
\begin{tabular}{@{}l l c c c c c c c c c c@{}}
\toprule
\multicolumn{12}{@{}l}{\textbf{a}\quad Author weights (E1)}\\
\addlinespace[2pt]
 & & \multicolumn{3}{c}{\colh{mAUPR, author metric}} & \multicolumn{5}{c}{\colh{AUPR, this protocol}} & & \\
\cmidrule(lr){3-5}\cmidrule(lr){6-10}
\colh{Model} & \colh{Source} & Paper & Function & Stored & \colh{mAUPR} & MA & HE & EX & SE & \colh{$\Delta$ metric} & \colh{Indep.\ fine} \\
\midrule
M2MRF & IDRiD & .672 & .676 & .645 & .639 & \cellcolor[HTML]{C0E4D8}.408 & \cellcolor[HTML]{88CEC9}.653 & \cellcolor[HTML]{5BB9C0}.832 & \cellcolor[HTML]{85CDC8}.664 & \down{.036} & \ph{.240} \\
M2MRF & DDR & .489 & .492 & .426 & .401 & \na & \na & \na & \na & \down{.091} & \ph{.276} \\
\addlinespace[2pt]
HRDecoder & IDRiD & \ph{.713} & \ph{.713} & \ph{.687} & \ph{.667} & \cellcolor[HTML]{B9E1D5}\ph{.449} & \cellcolor[HTML]{7CC9C6}\ph{.696} & \cellcolor[HTML]{59B7BF}\ph{.849} & \cellcolor[HTML]{83CCC8}\ph{.673} & \down{\ph{.047}} & \ph{.269} \\
HRDecoder & DDR & \ph{.492} & \ph{.492} & \ph{.449} & \ph{.429} & \cellcolor[HTML]{DFF0E5}\ph{.240} & \cellcolor[HTML]{AFDDD1}\ph{.500} & \cellcolor[HTML]{9CD6CC}\ph{.577} & \cellcolor[HTML]{C2E5D8}\ph{.400} & \down{\ph{.063}} & \ph{.310} \\
\addlinespace[2pt]
WSRFNet & IDRiD & \ph{.690} & \ph{.676} & \ph{.665} & \ph{.645} & \cellcolor[HTML]{BEE3D7}\ph{.424} & \cellcolor[HTML]{88CEC9}\ph{.650} & \cellcolor[HTML]{5AB8BF}\ph{.839} & \cellcolor[HTML]{85CDC8}\ph{.665} & \down{\ph{.032}} & \ph{.228} \\
WSRFNet & DDR & \ph{.500} & \ph{.494} & \ph{.425} & \ph{.405} & \cellcolor[HTML]{E1F1E6}\ph{.225} & \cellcolor[HTML]{B7E0D4}\ph{.465} & \cellcolor[HTML]{9ED7CD}\ph{.569} & \cellcolor[HTML]{C8E7DB}\ph{.364} & \down{\ph{.089}} & \ph{.254} \\
\bottomrule
\end{tabular}

\vspace{6pt}
\begin{tabular}{@{}l l c c c c c@{}}
\toprule
\multicolumn{7}{@{}l}{\textbf{b}\quad Author recipe, retrained (E1r)}\\
\addlinespace[2pt]
\colh{Model} & \colh{Source} & \colh{Paper} & \colh{Retrained, author metric} & \colh{Seeds} & \colh{$\Delta$ to paper} & \colh{Retrained, this protocol} \\
\midrule
M2MRF & IDRiD & .672 & \ph{.612}\sd{\ph{.154}} & \ph{3} & \down{\ph{.060}} & \ph{.549}\sd{\ph{.198}} \\
\quad first run, batch 1 & IDRiD & .672 & .576 & 1 & \down{.096} & .493 \\
M2MRF & DDR & .489 & \ph{.454}\sd{\ph{.002}} & \ph{3} & \down{\ph{.036}} & \ph{.416}\sd{\ph{.008}} \\
HACDR-Net & IDRiD & \ph{.700} & \ph{.596}\sd{\ph{.011}} & \ph{3} & \down{\ph{.104}} & \ph{.563}\sd{\ph{.016}} \\
HACDR-Net & DDR & \ph{.520} & \ph{.426}\sd{\ph{.006}} & \ph{3} & \down{\ph{.094}} & \ph{.389}\sd{\ph{.004}} \\
\bottomrule
\end{tabular}
\begin{tablenotes}[flushleft]\scriptsize
\item ``$\Delta$ metric'' is this protocol minus the authors' function on the same checkpoint. ``Indep.\ fine'' is the mean mAUPR on TJDR standard-field and ultra-widefield test sets. The M2MRF paper reports the mean of three runs (67.24 on IDRiD, 48.94 on DDR), the repository the best run (67.55, 49.20). HRDecoder: the released checkpoints (repository table 71.34, 49.24) through the authors' test pipeline (resize to $2880\times1920$ with sliding $1920$ windows on IDRiD, $2048\times2048$ whole image on DDR). The first M2MRF rerun used an effective batch of 1 instead of the authors' 4.
\phnote
\end{tablenotes}
\end{threeparttable}
\end{adjustbox}
\end{table*}


% ============================================================== R4
\subsection{A unified protocol to retrain the models}
Source recipes cannot serve as a common ground for retraining. The recipes of the general-purpose networks were written for natural images or abdominal organs; moved to IDRiD unchanged they differ about 67-fold in the number of training images seen, and those that resize the whole image to a small fixed input discard a large part of the small lesions (Fig.~\ref{fig:scale}b). We therefore fixed everything that is not part of a method: data and splits, a canvas of 1,440 pixels across the field of view, the training amount counted in pixels, augmentation, selection of the checkpoint by validation \mAUPR, inference geometry and evaluation. Each model keeps its loss, optimizer, learning rate, weight decay, schedule shape and effective batch (Methods and Table~\ref{tab:recipes}). The training amount was frozen from the validation curves of two models; no test metric was used. We retrained \rn{n.main} models on IDRiD and on DDR, with three seeds for the seven networks that have a source recipe and for the self-configuring baseline, and one seed for the adapted foundation models (two of them, which also serve as representative models below, with three seeds where noted); the rank tests use seed 0 of every model. The source recipes, for the models for which they completed, are reported in Supplementary Table~4.

\subsection{Performance of retrained models on internal tests}
\begin{simtext}
On the internal tests (Fig.~\ref{fig:retrained}a and Table~\ref{tab:main}), \rn{bstd.IDRiD.best.model} reached the highest \mAUPR{} on IDRiD (\rn{bstd.IDRiD.best.value}), followed by \rn{bstd.IDRiD.second.model} (\rn{bstd.IDRiD.second.value}); on DDR the first two were \rn{bstd.DDR.best.model} (\rn{bstd.DDR.best.value}) and \rn{bstd.DDR.second.model} (\rn{bstd.DDR.second.value}). The difference between the first two models was \rn{bstd.IDRiD.gap12} on IDRiD and \rn{bstd.DDR.gap12} on DDR, against a median seed standard deviation of \rn{bstd.IDRiD.sd.median} and \rn{bstd.DDR.sd.median}. IDRiD has 27 test images, and its bootstrap intervals are correspondingly wide; we use it to connect to the literature and base comparisons between architectures on DDR and on the independent tests.

Lesion classes differed more than models. On DDR the median AUPR over models was \rn{bstd.DDR.EX.median} for hard exudates and \rn{bstd.DDR.MA.median} for microaneurysms, and the best model reached \rn{bstd.DDR.MA.best} for microaneurysms. Networks with a coarse output stride or 16-pixel tokens lost most on this class. One failure was specific to a loss rather than to an architecture: the reference U-Net, trained on patches with a per-image Dice term smoothed by 1, drove its rare-class channels towards zero (soft exudate AUPR \rn{bstd.DDR.unet.SE} on DDR), because 83\% of DDR training patches contain no soft exudate and an empty patch rewards an all-zero output (Methods).
\end{simtext}

\begin{table*}[t]
\centering\tablefont
\caption{\textbf{Unified retraining (B-std): internal tests.} Every model is trained on one source under the shared protocol and scored on that source's official test set. Per-class values are exact pooled AUPR inside the field of view, mean of three seeds; mAUPR is the mean of the four classes $\pm$ s.d.\ over seeds; Dice is pooled, at the per-class threshold chosen on the validation set. Cell shading follows AUPR. \best{Bold}, best in column; \second{underlined}, second.}
\label{tab:main}
\begin{adjustbox}{max width=\textwidth}
\begin{threeparttable}
\begin{tabular}{@{}l r cccc c c @{\hspace{10pt}} cccc c c@{}}
\toprule
 & & \multicolumn{6}{c}{\colh{Trained and tested on IDRiD} {\scriptsize(44 / 10 / 27)}} & \multicolumn{6}{c}{\colh{Trained and tested on DDR} {\scriptsize(383 / 149 / 225)}} \\
\cmidrule(lr){3-8}\cmidrule(l){9-14}
\colh{Model} & \colh{Params} & MA & HE & EX & SE & \colh{mAUPR} & \colh{Dice} & MA & HE & EX & SE & \colh{mAUPR} & \colh{Dice} \\
\midrule
\grouprow{14}{\famchip{general}\; General segmentation}
U-Net & 7.8 & \cellcolor[HTML]{CAE8DC}\ph{.351} & \cellcolor[HTML]{99D5CC}\ph{.588} & \cellcolor[HTML]{66C1C2}\ph{.779} & \cellcolor[HTML]{E1F1E6}\ph{.222} & \ph{.485}\sd{\ph{.013}} & \ph{.478} & \cellcolor[HTML]{EEF5ED}\ph{.080} & \cellcolor[HTML]{E1F1E6}\ph{.225} & \cellcolor[HTML]{AFDDD1}\ph{.501} & \cellcolor[HTML]{EEF5ED}\ph{.080} & \ph{.222}\sd{\ph{.009}} & \ph{.271} \\
DeepLabv3 & \ph{60.9} & \cellcolor[HTML]{DEEFE4}\ph{.244} & \cellcolor[HTML]{A5D9CE}\ph{.539} & \cellcolor[HTML]{6FC4C4}\ph{.748} & \cellcolor[HTML]{A1D8CE}\ph{.556} & \ph{.522}\sd{\ph{.023}} & \ph{.506} & \cellcolor[HTML]{ECF4EB}\ph{.107} & \cellcolor[HTML]{CAE8DC}\ph{.356} & \cellcolor[HTML]{B8E1D4}\ph{.458} & \cellcolor[HTML]{D9EDE2}\ph{.270} & \ph{.298}\sd{\ph{.001}} & \ph{.328} \\
HRNet-W48 & \ph{65.8} & \cellcolor[HTML]{BBE2D6}\ph{.436} & \cellcolor[HTML]{88CEC9}\ph{.652} & \cellcolor[HTML]{5BB9C0}\ph{.833} & \cellcolor[HTML]{85CDC8}\ph{.663} & \ph{.646}\sd{\ph{.011}} & \ph{.603} & \cellcolor[HTML]{E2F1E6}\ph{.219} & \cellcolor[HTML]{B5E0D3}\ph{.469} & \cellcolor[HTML]{9CD6CC}\ph{.574} & \cellcolor[HTML]{C6E6DA}\ph{.373} & \ph{.409}\sd{\ph{.013}} & \ph{.416} \\
Swin-Unet & \ph{27.2} & \cellcolor[HTML]{DEEFE4}\ph{.242} & \cellcolor[HTML]{ACDCD0}\ph{.517} & \cellcolor[HTML]{71C5C4}\ph{.739} & \cellcolor[HTML]{A5D9CE}\ph{.543} & \ph{.511}\sd{\ph{.019}} & \ph{.500} & \cellcolor[HTML]{ECF5EC}\ph{.102} & \cellcolor[HTML]{D0EADE}\ph{.321} & \cellcolor[HTML]{C0E4D8}\ph{.407} & \cellcolor[HTML]{DCEFE4}\ph{.250} & \ph{.270}\sd{\ph{.006}} & \ph{.308} \\
FCT & \ph{51.4} & \cellcolor[HTML]{DCEFE4}\ph{.251} & \cellcolor[HTML]{AEDDD0}\ph{.508} & \cellcolor[HTML]{77C7C5}\ph{.721} & \cellcolor[HTML]{ACDCD0}\ph{.516} & \ph{.499}\sd{\ph{.012}} & \ph{.485} & \cellcolor[HTML]{ECF5EC}\ph{.105} & \cellcolor[HTML]{D3EBE0}\ph{.307} & \cellcolor[HTML]{C2E5D8}\ph{.400} & \cellcolor[HTML]{E0F0E5}\ph{.232} & \ph{.261}\sd{\ph{.015}} & \ph{.303} \\
\grouprow{14}{\famchip{auto}\; Self-configuring}
nnU-Net & \ph{33.4} & \cellcolor[HTML]{B3DFD2}\ph{\best{.479}} & \cellcolor[HTML]{81CBC7}\ph{\second{.679}} & \cellcolor[HTML]{58B6BF}\ph{\best{.856}} & \cellcolor[HTML]{7CC9C6}\ph{\best{.698}} & \ph{\best{.678}}\sd{\ph{.020}} & \ph{\second{.629}} & \cellcolor[HTML]{DEEFE4}\ph{\second{.248}} & \cellcolor[HTML]{B4DFD3}\ph{\second{.477}} & \cellcolor[HTML]{9CD6CC}\ph{\second{.579}} & \cellcolor[HTML]{C3E5D9}\ph{\second{.397}} & \ph{\second{.425}}\sd{\ph{.015}} & \ph{\second{.430}} \\
nnU-Net ResEnc-M & \ph{57.5} & \cellcolor[HTML]{B4DFD3}\ph{\second{.473}} & \cellcolor[HTML]{80CBC7}\ph{\best{.688}} & \cellcolor[HTML]{58B6BF}\ph{\second{.854}} & \cellcolor[HTML]{7ECAC6}\ph{\second{.692}} & \ph{\second{.677}}\sd{\ph{.006}} & \ph{\best{.629}} & \cellcolor[HTML]{DBEEE3}\ph{\best{.256}} & \cellcolor[HTML]{B0DED1}\ph{\best{.497}} & \cellcolor[HTML]{97D4CB}\ph{\best{.597}} & \cellcolor[HTML]{C2E5D8}\ph{\best{.402}} & \ph{\best{.438}}\sd{\ph{.022}} & \ph{\best{.448}} \\
\grouprow{14}{\famchip{fundus}\; Fundus-specific}
M2MRF & 70.3 & \cellcolor[HTML]{BDE3D6}\ph{.427} & \cellcolor[HTML]{8ED0CA}\ph{.628} & \cellcolor[HTML]{5DBAC0}\ph{.823} & \cellcolor[HTML]{88CEC9}\ph{.649} & \ph{.632}\sd{\ph{.007}} & \ph{.594} & \cellcolor[HTML]{E0F0E5}\ph{.234} & \cellcolor[HTML]{B5E0D3}\ph{.470} & \cellcolor[HTML]{9ED7CD}\ph{.568} & \cellcolor[HTML]{C5E6DA}\ph{.380} & \ph{.413}\sd{\ph{.022}} & \ph{.420} \\
HACDR-Net & \ph{49.6} & \cellcolor[HTML]{CFEADE}\ph{.323} & \cellcolor[HTML]{9CD6CC}\ph{.577} & \cellcolor[HTML]{65C0C2}\ph{.783} & \cellcolor[HTML]{99D5CC}\ph{.591} & \ph{.569}\sd{\ph{.019}} & \ph{.543} & \cellcolor[HTML]{E6F3E8}\ph{.169} & \cellcolor[HTML]{C0E4D8}\ph{.407} & \cellcolor[HTML]{AEDDD0}\ph{.511} & \cellcolor[HTML]{D0EADE}\ph{.320} & \ph{.352}\sd{\ph{.017}} & \ph{.377} \\
\grouprow{14}{\famchip{foundation}\; Foundation models, adapted}
SAMed & \ph{91.2} & \cellcolor[HTML]{DEEFE4}\ph{.243} & \cellcolor[HTML]{A5D9CE}\ph{.544} & \cellcolor[HTML]{6AC2C3}\ph{.766} & \cellcolor[HTML]{A0D7CD}\ph{.565} & \ph{.530}\sd{\ph{.009}} & \ph{.512} & \cellcolor[HTML]{EDF5EC}\ph{.096} & \cellcolor[HTML]{CDE9DD}\ph{.340} & \cellcolor[HTML]{BAE2D5}\ph{.445} & \cellcolor[HTML]{DAEEE3}\ph{.267} & \ph{.287}\sd{\ph{.004}} & \ph{.327} \\
SAM-Adapter & 641.3 & \cellcolor[HTML]{E2F1E6}\ph{.216} & \cellcolor[HTML]{AADCCF}\ph{.523} & \cellcolor[HTML]{75C6C5}\ph{.724} & \cellcolor[HTML]{A7DACF}\ph{.533} & \ph{.499}\sd{\ph{.019}} & \ph{.491} & \cellcolor[HTML]{EEF5ED}\ph{.085} & \cellcolor[HTML]{D0EADE}\ph{.322} & \cellcolor[HTML]{C2E5D8}\ph{.401} & \cellcolor[HTML]{DFF0E5}\ph{.239} & \ph{.262}\sd{\ph{.009}} & \ph{.306} \\
SAM2-UNet & 216.5 & \cellcolor[HTML]{C9E7DB}\ph{.361} & \cellcolor[HTML]{8ED0CA}\ph{.632} & \cellcolor[HTML]{5EBBC0}\ph{.817} & \cellcolor[HTML]{8CD0C9}\ph{.640} & \ph{.613}\sd{\ph{.026}} & \ph{.578} & \cellcolor[HTML]{E5F2E7}\ph{.194} & \cellcolor[HTML]{BAE2D5}\ph{.442} & \cellcolor[HTML]{A3D9CE}\ph{.551} & \cellcolor[HTML]{C8E7DB}\ph{.366} & \ph{.388}\sd{\ph{.010}} & \ph{.403} \\
MedSAM (full FT) & 93.7 & \cellcolor[HTML]{E3F1E6}\ph{.211} & \cellcolor[HTML]{A7DACF}\ph{.535} & \cellcolor[HTML]{6CC3C3}\ph{.757} & \cellcolor[HTML]{A5D9CE}\ph{.540} & \ph{.511}\sd{\ph{.022}} & \ph{.493} & \cellcolor[HTML]{EDF5EC}\ph{.089} & \cellcolor[HTML]{CEE9DD}\ph{.330} & \cellcolor[HTML]{BDE3D6}\ph{.428} & \cellcolor[HTML]{DBEEE3}\ph{.257} & \ph{.276}\sd{\ph{.001}} & \ph{.311} \\
RETFound + UPerNet & 352.6 & \cellcolor[HTML]{D4ECE0}\ph{.302} & \cellcolor[HTML]{95D3CB}\ph{.604} & \cellcolor[HTML]{63BFC1}\ph{.795} & \cellcolor[HTML]{91D2CA}\ph{.616} & \ph{.579}\sd{\ph{.013}} & \ph{.555} & \cellcolor[HTML]{E9F4EA}\ph{.134} & \cellcolor[HTML]{C3E5D9}\ph{.394} & \cellcolor[HTML]{AFDDD1}\ph{.499} & \cellcolor[HTML]{D1EBDF}\ph{.311} & \ph{.334}\sd{\ph{.024}} & \ph{.358} \\
\bottomrule
\end{tabular}
\begin{tablenotes}[flushleft]\scriptsize
\item Params in millions. nnU-Net uses its own preprocessing and planning with the same splits, test images and evaluation. H2Former is registered in tier~1 but deferred: at its 960-pixel input one training step exceeds the 40\,GB of the cards used here.
\phnote
\end{tablenotes}
\end{threeparttable}
\end{adjustbox}
\end{table*}


\subsection{Performance of retrained models on independent tests}
\begin{simtext}
Applied to sources that were never used for training, every model lost performance (Fig.~\ref{fig:retrained}b,c and Table~\ref{tab:independent}). For models trained on DDR the mean loss in \mAUPR{} was \rn{indep.DDR.TJDR-std.drop} on the standard-field TJDR test set (Holm-adjusted $P$ \rn{indep.DDR.TJDR-std.p}), \rn{indep.DDR.TJDR-uwf.drop} on the ultra-widefield set, \rn{indep.DDR.FGADR.drop} on FGADR ($P$ \rn{indep.DDR.FGADR.p}) and \rn{indep.DDR.Retinal-Lesions.drop} on Retinal-Lesions. For models trained on the 44 IDRiD images the losses were larger (\rn{indep.IDRiD.TJDR-std.drop}, \rn{indep.IDRiD.TJDR-uwf.drop}, \rn{indep.IDRiD.FGADR.drop} and \rn{indep.IDRiD.Retinal-Lesions.drop}). The order of the sources was the same for both: least on fine labels at a standard field of view, more on mixed and coarse labels, and most where the imaging geometry changes.

Internal and independent performance were correlated across models (Spearman $\rho$ = \rn{indep.DDR.rho} for DDR), and the model that led on the fine-grained independent tests (\rn{indep.DDR.top.model}, \rn{indep.DDR.top.value}) was \rn{indep.DDR.sametop} as on the internal test. Averaged by family, the self-configuring baseline was ahead on internal (\rn{family.auto.internal}) and on fine-grained independent tests (\rn{family.auto.independent}), fundus-specific networks (\rn{family.fundus.internal} and \rn{family.fundus.independent}) were close to the best general network, and adapted foundation models did not exceed either (Fig.~\ref{fig:retrained}d).

Over the seven test sets the Friedman test rejected equal ranks ($\chi^2$ = \rn{cd.chi2}, $P$ \rn{cd.p}). \rn{cd.first.model} had the best average rank (\rn{cd.first.rank}) and \rn{cd.last.model} the worst (\rn{cd.last.rank}); with seven test sets the Nemenyi critical difference is \rn{cd.value} ranks, and \rn{cd.ntop} of \rn{n.main} models were not separable from the first (Fig.~\ref{fig:retrained}e). The ranking published for the \rn{tau.published.n} models that report IDRiD results agreed only partly with the ranking after retraining (Kendall $\tau$ = \rn{tau.published}, $P$ = \rn{tau.published.p}; Fig.~\ref{fig:retrained}f). Using the source recipes instead of the unified protocol lowered \rn{b1.IDRiD.nlower} of \rn{b1.n} models on IDRiD, by up to \rn{b1.IDRiD.worst.delta} for \rn{b1.IDRiD.worst.model} (Supplementary Table~4), and a threefold change of the learning rate in either direction, tested on DDR for one general and one fundus-specific network, did not change their order (Extended Data Fig.~7).
\end{simtext}

\begin{figure}[p]
\centering
\includegraphics[width=\textwidth,height=0.78\textheight,keepaspectratio]{\figfile{fig4_retrained}}
\caption{\textbf{Performance of retrained models under the unified protocol.} \textbf{a}, Internal tests: per-class AUPR (mean of three seeds) and \mAUPR{} (bar; error bar, s.d.\ over seeds). Models are ordered by their average rank in \textbf{e}; the square gives the family. \textbf{b}, \mAUPR{} of every model on the seven test sets, by training source. Int., internal test; Cross, test set of the other training source; the bar under each column gives the annotation granularity. Columns on the right: internal test, mean of the two fine-grained independent test sets, and their difference. \textbf{c}, The same models on the internal test and on four independent test sets; each line is one model. $P$ values, two-sided Wilcoxon signed-rank tests across models with Holm correction ($n$ = \rn{n.main}). \textbf{d}, \mAUPR{} by family on internal and on fine-grained independent tests; coloured dots, one model and training source; black, mean $\pm$ s.d.; K--W, Kruskal--Wallis test. \textbf{e}, Critical-difference diagram over the seven test sets; bars join models whose average ranks differ by less than the Nemenyi critical difference. \textbf{f}, Rank on IDRiD as published and after retraining.}
\label{fig:retrained}
\end{figure}

\begin{table*}[t]
\centering\tablefont\setlength{\tabcolsep}{3.3pt}
\caption{\textbf{Unified retraining (B-std): independent tests.} The checkpoints of Table~\ref{tab:main} are applied without any adaptation to five test sets from sources never used for training. Values are mAUPR (mean of three seeds); the DiaRetDB1 column is the official image-level ROC AUC, averaged over the four lesion types, because its region-consensus ground truth biases pixel precision. $\Delta$ is the mean of the two fine-grained TJDR test sets minus the internal test. Average rank is over the seven test sets of Fig.~\ref{fig:retrained}e (Friedman $\chi^2$ = \ph{58.1}; Nemenyi critical difference \ph{7.50}).}
\label{tab:independent}
\begin{adjustbox}{max width=\textwidth}
\begin{threeparttable}
\begin{tabular}{@{}l c cccc c c @{\hspace{8pt}} c cccc c c @{\hspace{8pt}} c@{}}
\toprule
 & \multicolumn{7}{c}{\colh{Trained on IDRiD}} & \multicolumn{7}{c}{\colh{Trained on DDR}} & \\
\cmidrule(lr){2-8}\cmidrule(lr){9-15}
 & & \multicolumn{2}{c}{TJDR} & & & & & & \multicolumn{2}{c}{TJDR} & & & & & \colh{Avg.} \\
\colh{Model} & Int. & Std & UWF & FGADR & Ret.-L. & DiaRet & $\Delta$ & Int. & Std & UWF & FGADR & Ret.-L. & DiaRet & $\Delta$ & \colh{rank} \\
 & \tier{fine} & \tier{fine} & \tier{fine} & \tier{mixed} & \tier{coarse} & \tier{region} &  & \tier{fine} & \tier{fine} & \tier{fine} & \tier{mixed} & \tier{coarse} & \tier{region} &  & \\
\midrule
\grouprow{16}{\famchip{general}\; General segmentation}
U-Net & \ph{.66} & \cellcolor[HTML]{D4ECE0}\ph{.23} & \cellcolor[HTML]{EBF4EB}\ph{.10} & \cellcolor[HTML]{E6F3E8}\ph{.13} & \cellcolor[HTML]{E9F3EA}\ph{.11} & \ph{.77} & \cellcolor[HTML]{E47F75}\ph{$$-$$.49} & \ph{.22} & \cellcolor[HTML]{D9EDE2}\ph{.21} & \cellcolor[HTML]{EBF4EB}\ph{.09} & \cellcolor[HTML]{E8F3E9}\ph{.12} & \cellcolor[HTML]{EBF4EB}\ph{.10} & \ph{.77} & \cellcolor[HTML]{EFDFDB}\ph{$$-$$.07} & \ph{12.29} \\
DeepLabv3 & \ph{.54} & \cellcolor[HTML]{CAE8DC}\ph{.27} & \cellcolor[HTML]{E4F2E7}\ph{.16} & \cellcolor[HTML]{E1F1E6}\ph{.18} & \cellcolor[HTML]{E7F3E9}\ph{.13} & \ph{.78} & \cellcolor[HTML]{EAA69E}\ph{$$-$$.32} & \ph{.38} & \cellcolor[HTML]{C3E5D9}\ph{.30} & \cellcolor[HTML]{E2F1E6}\ph{.17} & \cellcolor[HTML]{DCEFE4}\ph{.20} & \cellcolor[HTML]{E5F2E8}\ph{.14} & \ph{.79} & \cellcolor[HTML]{EECFCA}\ph{$$-$$.14} & \ph{8.29} \\
HRNet-W48 & \ph{\best{.70}} & \cellcolor[HTML]{CEE9DD}\ph{.26} & \cellcolor[HTML]{E4F2E7}\ph{.16} & \cellcolor[HTML]{D4ECE0}\ph{.23} & \cellcolor[HTML]{E8F3E9}\ph{.12} & \ph{\second{.86}} & \cellcolor[HTML]{E47F75}\ph{$$-$$.49} & \ph{\best{.45}} & \cellcolor[HTML]{C4E5D9}\ph{.30} & \cellcolor[HTML]{E0F0E5}\ph{.18} & \cellcolor[HTML]{CBE8DC}\ph{.27} & \cellcolor[HTML]{E6F2E8}\ph{.14} & \ph{\second{.87}} & \cellcolor[HTML]{ECC0BA}\ph{$$-$$.21} & \ph{4.43} \\
Swin-Unet & \ph{\second{.69}} & \cellcolor[HTML]{D3EBE0}\ph{.24} & \cellcolor[HTML]{EBF4EB}\ph{.10} & \cellcolor[HTML]{E3F2E7}\ph{.16} & \cellcolor[HTML]{E9F3EA}\ph{.11} & \ph{.76} & \cellcolor[HTML]{E37D74}\ph{$$-$$.52} & \ph{.43} & \cellcolor[HTML]{CBE8DC}\ph{.27} & \cellcolor[HTML]{E9F3EA}\ph{.11} & \cellcolor[HTML]{DFF0E5}\ph{.18} & \cellcolor[HTML]{E8F3E9}\ph{.12} & \ph{.78} & \cellcolor[HTML]{ECB8B1}\ph{$$-$$.24} & \ph{9.43} \\
FCT & \ph{.30} & \cellcolor[HTML]{CBE8DC}\ph{.27} & \cellcolor[HTML]{EAF4EA}\ph{.10} & \cellcolor[HTML]{E5F2E8}\ph{.14} & \cellcolor[HTML]{EAF4EA}\ph{.10} & \ph{.78} & \cellcolor[HTML]{EED5D1}\ph{$$-$$.11} & \ph{.18} & \cellcolor[HTML]{C6E6DA}\ph{.29} & \cellcolor[HTML]{E9F3EA}\ph{.11} & \cellcolor[HTML]{E3F2E7}\ph{.16} & \cellcolor[HTML]{E9F3EA}\ph{.11} & \ph{.79} & \cellcolor[HTML]{EAECEC}\ph{+.03} & \ph{12.14} \\
\grouprow{16}{\famchip{auto}\; Self-configuring}
nnU-Net & \ph{.67} & \cellcolor[HTML]{C3E5D9}\ph{.31} & \cellcolor[HTML]{DFF0E5}\ph{\best{.19}} & \cellcolor[HTML]{D0EADE}\ph{\second{.25}} & \cellcolor[HTML]{D8EDE2}\ph{\best{.22}} & \ph{.85} & \cellcolor[HTML]{E98F85}\ph{$$-$$.42} & \ph{\second{.43}} & \cellcolor[HTML]{BEE3D7}\ph{.33} & \cellcolor[HTML]{DBEEE3}\ph{.20} & \cellcolor[HTML]{CBE8DC}\ph{.27} & \cellcolor[HTML]{D5ECE1}\ph{\best{.22}} & \ph{.86} & \cellcolor[HTML]{EDC8C3}\ph{$$-$$.17} & \ph{\best{2.86}} \\
nnU-Net ResEnc-M & \ph{.63} & \cellcolor[HTML]{AEDDD0}\ph{\best{.39}} & \cellcolor[HTML]{E4F2E7}\ph{.15} & \cellcolor[HTML]{C5E6DA}\ph{\best{.29}} & \cellcolor[HTML]{E3F2E7}\ph{.16} & \ph{\best{.86}} & \cellcolor[HTML]{EA9D95}\ph{$$-$$.36} & \ph{.43} & \cellcolor[HTML]{A1D8CE}\ph{\best{.43}} & \cellcolor[HTML]{E0F0E5}\ph{.18} & \cellcolor[HTML]{BEE3D7}\ph{\best{.33}} & \cellcolor[HTML]{E2F1E6}\ph{.17} & \ph{\best{.88}} & \cellcolor[HTML]{EED4CF}\ph{$$-$$.12} & \ph{\second{3.29}} \\
\grouprow{16}{\famchip{fundus}\; Fundus-specific}
M2MRF & \ph{.66} & \cellcolor[HTML]{C4E5D9}\ph{.30} & \cellcolor[HTML]{E3F1E6}\ph{.16} & \cellcolor[HTML]{D4ECE0}\ph{.23} & \cellcolor[HTML]{E5F2E8}\ph{.14} & \ph{.82} & \cellcolor[HTML]{E98F85}\ph{$$-$$.43} & \ph{.37} & \cellcolor[HTML]{B7E0D4}\ph{.36} & \cellcolor[HTML]{DAEEE3}\ph{\second{.21}} & \cellcolor[HTML]{C8E7DB}\ph{\second{.28}} & \cellcolor[HTML]{E2F1E6}\ph{.17} & \ph{.84} & \cellcolor[HTML]{EFDCD8}\ph{$$-$$.09} & \ph{4.57} \\
HACDR-Net & \ph{.56} & \cellcolor[HTML]{C8E7DB}\ph{.29} & \cellcolor[HTML]{E8F3E9}\ph{.12} & \cellcolor[HTML]{E4F2E7}\ph{.15} & \cellcolor[HTML]{E8F3E9}\ph{.12} & \ph{.79} & \cellcolor[HTML]{EA9D95}\ph{$$-$$.36} & \ph{.27} & \cellcolor[HTML]{BAE2D5}\ph{.34} & \cellcolor[HTML]{E4F2E7}\ph{.15} & \cellcolor[HTML]{DCEFE4}\ph{.20} & \cellcolor[HTML]{E5F2E8}\ph{.14} & \ph{.82} & \cellcolor[HTML]{F0EBE8}\ph{$$-$$.02} & \ph{8.86} \\
\grouprow{16}{\famchip{foundation}\; Foundation models, adapted}
SAMed & \ph{.53} & \cellcolor[HTML]{D3EBE0}\ph{.24} & \cellcolor[HTML]{E7F3E9}\ph{.13} & \cellcolor[HTML]{DCEFE4}\ph{.20} & \cellcolor[HTML]{E5F2E8}\ph{.14} & \ph{.81} & \cellcolor[HTML]{EAA199}\ph{$$-$$.34} & \ph{.29} & \cellcolor[HTML]{CEE9DD}\ph{.26} & \cellcolor[HTML]{E5F2E8}\ph{.15} & \cellcolor[HTML]{D9EDE2}\ph{.21} & \cellcolor[HTML]{E5F2E7}\ph{.15} & \ph{.84} & \cellcolor[HTML]{EFDCD8}\ph{$$-$$.08} & \ph{8.57} \\
SAM-Adapter & \ph{.50} & \cellcolor[HTML]{D8EDE2}\ph{.21} & \cellcolor[HTML]{E7F3E9}\ph{.13} & \cellcolor[HTML]{E1F1E6}\ph{.18} & \cellcolor[HTML]{EAF4EA}\ph{.10} & \ph{.80} & \cellcolor[HTML]{EAA49C}\ph{$$-$$.33} & \ph{.26} & \cellcolor[HTML]{D4ECE0}\ph{.23} & \cellcolor[HTML]{E5F2E8}\ph{.14} & \cellcolor[HTML]{DCEFE4}\ph{.19} & \cellcolor[HTML]{E9F3EA}\ph{.11} & \ph{.80} & \cellcolor[HTML]{EFDEDA}\ph{$$-$$.07} & \ph{11.43} \\
SAM2-UNet & \ph{.61} & \cellcolor[HTML]{BBE2D6}\ph{\second{.34}} & \cellcolor[HTML]{E0F0E5}\ph{\second{.18}} & \cellcolor[HTML]{D9EDE2}\ph{.21} & \cellcolor[HTML]{E4F2E7}\ph{.15} & \ph{.84} & \cellcolor[HTML]{EA9F97}\ph{$$-$$.35} & \ph{.39} & \cellcolor[HTML]{AFDDD1}\ph{\second{.39}} & \cellcolor[HTML]{D8EDE2}\ph{\best{.22}} & \cellcolor[HTML]{CFEADE}\ph{.25} & \cellcolor[HTML]{E0F0E5}\ph{.18} & \ph{.87} & \cellcolor[HTML]{EFDCD8}\ph{$$-$$.08} & \ph{3.86} \\
MedSAM (full FT) & \ph{.51} & \cellcolor[HTML]{CDE9DD}\ph{.26} & \cellcolor[HTML]{E8F3E9}\ph{.12} & \cellcolor[HTML]{D6EDE1}\ph{.22} & \cellcolor[HTML]{E7F3E9}\ph{.13} & \ph{.80} & \cellcolor[HTML]{EAA69E}\ph{$$-$$.32} & \ph{.28} & \cellcolor[HTML]{C8E7DB}\ph{.28} & \cellcolor[HTML]{E6F3E8}\ph{.13} & \cellcolor[HTML]{D0EADE}\ph{.25} & \cellcolor[HTML]{E6F3E8}\ph{.13} & \ph{.81} & \cellcolor[HTML]{EFDFDB}\ph{$$-$$.07} & \ph{9.43} \\
RETFound + UPerNet & \ph{.58} & \cellcolor[HTML]{C2E5D8}\ph{.31} & \cellcolor[HTML]{E3F2E7}\ph{.16} & \cellcolor[HTML]{E0F0E5}\ph{.18} & \cellcolor[HTML]{DCEFE4}\ph{\second{.20}} & \ph{.84} & \cellcolor[HTML]{EAA199}\ph{$$-$$.35} & \ph{.33} & \cellcolor[HTML]{BDE3D6}\ph{.33} & \cellcolor[HTML]{E1F1E6}\ph{.18} & \cellcolor[HTML]{DBEEE3}\ph{.20} & \cellcolor[HTML]{DAEEE3}\ph{\second{.21}} & \ph{.86} & \cellcolor[HTML]{EFDCD8}\ph{$$-$$.08} & \ph{5.57} \\
\bottomrule
\end{tabular}
\begin{tablenotes}[flushleft]\scriptsize
\item Granularity of the test set: \tier{fine} fine \quad \tier{mixed} mixed \quad \tier{coarse} coarse \quad \tier{region} region consensus. Int., internal test (official test set of the training source). Std, standard field (35--50$^\circ$); UWF, ultra-widefield (133$^\circ$), predicted at native resolution with sliding windows. Retinal-Lesions excludes its grey-level 127 regions.
\phnote
\end{tablenotes}
\end{threeparttable}
\end{adjustbox}
\end{table*}


% ============================================================== R5
\subsection{Foundation models under four kinds of use}
Results for segmentation foundation models are often quoted without saying how the model was used. We separated four kinds of use and report them side by side without ranking one against another (Fig.~\ref{fig:foundation}a and Table~\ref{tab:foundation}): zero-shot use without any label; prompts derived from the ground truth, which is an upper bound and not a usable setting; parameter-efficient fine-tuning; and an encoder with a trained decoder or full fine-tuning.

\begin{simtext}
Zero-shot use did not segment DR lesions. Automatic masks of SAM and SAM~2, even when each mask was given the class it overlaps most, and text prompts of SAM~3 stayed far below every trained model (best: \rn{fm.zero.model}, \mAUPR{} \rn{fm.zero.value} on IDRiD). Text prompts were sensitive to wording: across three phrasings per lesion the best and worst AUPR differed \rn{fm.wording.ratio}-fold on average (Fig.~\ref{fig:foundation}c). With a box around every ground-truth component SAM~2 reached \rn{fm.oracle.all}, but an IDRiD image carries on average 44 microaneurysm and 142 hard exudate components; with 5 and 20 boxes per image and class the \mAUPR{} was \rn{fm.oracle.five} and \rn{fm.oracle.twenty} (Fig.~\ref{fig:foundation}b). Given the same boxes, the supervised models drew the lesions at least as well inside the box, so the gap of the promptable models lies in finding lesions, not in delineating them.

Among the trained settings, \rn{fm.peft.model} was the best parameter-efficient adaptation (\rn{fm.peft.value} on DDR, against \rn{fm.ref.nnunet} for nnU-Net). Under one UPerNet head, one recipe and one training amount on DDR, an encoder pretrained on retinal images (\rn{fm.enc.retfound.int} internal, \rn{fm.enc.retfound.ind} independent) was compared with its own starting point, a masked-autoencoder trained on ImageNet (\rn{fm.enc.imagenet.int} and \rn{fm.enc.imagenet.ind}), and the same contrast was repeated under the DINOv2 objective (retinal \rn{fm.enc.retfound.ind} against general \rn{fm.enc.dinov2.ind} on independent data), so that the effect of retinal pretraining does not depend on one pretraining objective (Fig.~\ref{fig:foundation}d,e).
\end{simtext}

\begin{figure}[t]
\centering
\widefig{fig5_foundation}
\caption{\textbf{Foundation models under four kinds of use.} \textbf{a}, Per-class AUPR and \mAUPR{} on the IDRiD and DDR test sets. Zero-shot and oracle rows use no training image; the two trained settings are trained on the source of the test set. \textbf{b}, Oracle prompts: AUPR on IDRiD against the number of ground-truth boxes or points given per image and class; dotted line, nnU-Net trained on IDRiD. \textbf{c}, SAM~3 with three wordings per lesion class. \textbf{d}, Four ViT-L encoders under the same UPerNet head and recipe, trained on DDR (seed 0): internal test and mean of the two TJDR test sets; lines join each retinal encoder to the natural-image encoder trained with the same objective. \textbf{e}, Best model of each kind of use against the fundus-specific and self-configuring references.}
\label{fig:foundation}
\end{figure}

\begin{table*}[t]
\centering\tablefont
\caption{\textbf{Foundation models under four kinds of use.} The kinds are reported in separate blocks and are not ranked against each other: zero-shot use sees no test label (its density, area cap and wording are chosen on the validation pool); the upper-bound block reads the test labels (masks classed by their overlap with the ground truth, prompts derived from every ground-truth component) and is not a usable setting; the two trained settings follow the unified protocol. AUPR on the IDRiD and DDR test sets. \best{Bold} and \second{underlined} mark the best and second-best rows of each column within a block.}
\label{tab:foundation}
\begin{adjustbox}{max width=\textwidth}
\begin{threeparttable}
\begin{tabular}{@{}l l l cccc c @{\hspace{10pt}} cccc c@{}}
\toprule
 & & & \multicolumn{5}{c}{\colh{IDRiD test}} & \multicolumn{5}{c}{\colh{DDR test}} \\
\cmidrule(lr){4-8}\cmidrule(l){9-13}
\colh{Model} & \colh{Use} & \colh{Backbone} & MA & HE & EX & SE & \colh{mAUPR} & MA & HE & EX & SE & \colh{mAUPR} \\
\midrule
\grouprow{13}{Zero-shot: no test label; density, area cap and wording chosen on the validation pool (summary = any-lesion AUPR$^\ddagger$)}
\famchip{foundation}\; SAM & automatic masks, class-agnostic & ViT-H & \na & \na & \na & \na & \ph{.06} & \na & \na & \na & \na & \ph{.06} \\
\famchip{foundation}\; SAM 2 & automatic masks, class-agnostic & Hiera-L & \na & \na & \na & \na & \ph{\second{.07}} & \na & \na & \na & \na & \ph{\second{.07}} \\
\famchip{foundation}\; SAM 3 & text, wording chosen on validation & -- & \cellcolor[HTML]{F3F7EF}\ph{.03} & \cellcolor[HTML]{E5F2E7}\ph{.19} & \cellcolor[HTML]{D6EDE1}\ph{.29} & \cellcolor[HTML]{E2F1E6}\ph{.22} & \ph{\best{.16}} & \cellcolor[HTML]{F3F7EF}\ph{.03} & \cellcolor[HTML]{E9F4EA}\ph{.14} & \cellcolor[HTML]{DEEFE4}\ph{.25} & \cellcolor[HTML]{E3F1E6}\ph{.21} & \ph{\best{.12}} \\
\grouprow{13}{Upper bounds that read the test labels (not a usable setting)}
\famchip{foundation}\; SAM & automatic masks, class by overlap$^\dagger$ & ViT-H & \cellcolor[HTML]{F3F7F0}\ph{.02} & \cellcolor[HTML]{EBF4EB}\ph{.12} & \cellcolor[HTML]{E6F2E8}\ph{.18} & \cellcolor[HTML]{EAF4EA}\ph{.13} & \ph{.11} & \cellcolor[HTML]{F3F7F0}\ph{.02} & \cellcolor[HTML]{EBF4EB}\ph{.12} & \cellcolor[HTML]{E8F3E9}\ph{.15} & \cellcolor[HTML]{E9F4EA}\ph{.13} & \ph{.11} \\
\famchip{foundation}\; SAM 2 & automatic masks, class by overlap$^\dagger$ & Hiera-L & \cellcolor[HTML]{F3F7F0}\ph{.02} & \cellcolor[HTML]{E8F3E9}\ph{.15} & \cellcolor[HTML]{E8F3E9}\ph{.16} & \cellcolor[HTML]{E0F0E5}\ph{.23} & \ph{.14} & \cellcolor[HTML]{F2F7EF}\ph{.04} & \cellcolor[HTML]{E9F4EA}\ph{.14} & \cellcolor[HTML]{E2F1E6}\ph{.22} & \cellcolor[HTML]{DEEFE4}\ph{.25} & \ph{.16} \\
\famchip{foundation}\; SAM & oracle boxes, every component & ViT-B & \cellcolor[HTML]{D4ECE0}\ph{\second{.30}} & \cellcolor[HTML]{93D2CB}\ph{\second{.61}} & \cellcolor[HTML]{87CDC8}\ph{\second{.66}} & \cellcolor[HTML]{77C7C5}\ph{\second{.72}} & \ph{\second{.57}} & \cellcolor[HTML]{DBEEE3}\ph{\second{.26}} & \cellcolor[HTML]{A9DBCF}\ph{\second{.53}} & \cellcolor[HTML]{9AD5CC}\ph{\second{.58}} & \cellcolor[HTML]{8CD0C9}\ph{\second{.64}} & \ph{\second{.50}} \\
\famchip{foundation}\; SAM 2 & oracle points, every component & Hiera-L & \cellcolor[HTML]{E3F1E6}\ph{.21} & \cellcolor[HTML]{AFDDD1}\ph{.50} & \cellcolor[HTML]{ACDCD0}\ph{.52} & \cellcolor[HTML]{91D2CA}\ph{.62} & \ph{.46} & \cellcolor[HTML]{E5F2E8}\ph{.19} & \cellcolor[HTML]{B9E1D5}\ph{.45} & \cellcolor[HTML]{B8E1D4}\ph{.45} & \cellcolor[HTML]{A0D7CD}\ph{.56} & \ph{.41} \\
\famchip{foundation}\; SAM 2 & oracle boxes, every component & Hiera-L & \cellcolor[HTML]{CDE9DD}\ph{\best{.34}} & \cellcolor[HTML]{88CEC9}\ph{\best{.65}} & \cellcolor[HTML]{7CC9C6}\ph{\best{.70}} & \cellcolor[HTML]{6EC4C3}\ph{\best{.75}} & \ph{\best{.61}} & \cellcolor[HTML]{D6EDE1}\ph{\best{.29}} & \cellcolor[HTML]{99D5CC}\ph{\best{.59}} & \cellcolor[HTML]{93D2CB}\ph{\best{.61}} & \cellcolor[HTML]{85CDC8}\ph{\best{.67}} & \ph{\best{.54}} \\
\famchip{foundation}\; MedSAM & oracle boxes, every component & ViT-B & \cellcolor[HTML]{E3F2E7}\ph{.20} & \cellcolor[HTML]{9ED7CD}\ph{.57} & \cellcolor[HTML]{9CD6CC}\ph{.58} & \cellcolor[HTML]{7ECAC6}\ph{.69} & \ph{.51} & \cellcolor[HTML]{E4F2E7}\ph{.20} & \cellcolor[HTML]{AFDDD1}\ph{.51} & \cellcolor[HTML]{AEDDD0}\ph{.51} & \cellcolor[HTML]{90D1CA}\ph{.62} & \ph{.46} \\
\grouprow{13}{Parameter-efficient fine-tuning (trained on the source of the test set)}
\famchip{foundation}\; SAMed & LoRA & ViT-B & \cellcolor[HTML]{E1F1E6}\ph{\second{.23}} & \cellcolor[HTML]{A9DBCF}\ph{\second{.53}} & \cellcolor[HTML]{6EC4C3}\ph{\second{.75}} & \cellcolor[HTML]{A7DACF}\ph{\second{.54}} & \ph{\second{.51}} & \cellcolor[HTML]{ECF5EC}\ph{\second{.11}} & \cellcolor[HTML]{C8E7DB}\ph{\second{.37}} & \cellcolor[HTML]{B8E1D4}\ph{\second{.46}} & \cellcolor[HTML]{D8EDE2}\ph{\second{.28}} & \ph{\second{.30}} \\
\famchip{foundation}\; SAM-Adapter & adapters & ViT-H & \cellcolor[HTML]{E3F1E6}\ph{.21} & \cellcolor[HTML]{AEDDD0}\ph{.51} & \cellcolor[HTML]{73C6C4}\ph{.73} & \cellcolor[HTML]{AADCCF}\ph{.52} & \ph{.49} & \cellcolor[HTML]{EDF5EC}\ph{.09} & \cellcolor[HTML]{CFEADE}\ph{.33} & \cellcolor[HTML]{BEE3D7}\ph{.42} & \cellcolor[HTML]{DCEFE4}\ph{.25} & \ph{.27} \\
\famchip{foundation}\; SAM2-UNet & adapters + U-shaped decoder & Hiera-L & \cellcolor[HTML]{C6E6DA}\ph{\best{.38}} & \cellcolor[HTML]{8ACFC9}\ph{\best{.65}} & \cellcolor[HTML]{5DBAC0}\ph{\best{.82}} & \cellcolor[HTML]{88CEC9}\ph{\best{.65}} & \ph{\best{.62}} & \cellcolor[HTML]{E5F2E8}\ph{\best{.18}} & \cellcolor[HTML]{BAE2D5}\ph{\best{.44}} & \cellcolor[HTML]{A3D9CE}\ph{\best{.55}} & \cellcolor[HTML]{CAE8DC}\ph{\best{.35}} & \ph{\best{.38}} \\
\grouprow{13}{Encoder + decoder (trained on the source of the test set)}
\famchip{foundation}\; MedSAM & full fine-tuning & ViT-B & \cellcolor[HTML]{E2F1E6}\ph{\second{.22}} & \cellcolor[HTML]{A7DACF}\ph{\second{.54}} & \cellcolor[HTML]{6CC3C3}\ph{\second{.76}} & \cellcolor[HTML]{A1D8CE}\ph{\second{.56}} & \ph{\second{.52}} & \cellcolor[HTML]{EDF5EC}\ph{\second{.10}} & \cellcolor[HTML]{CBE8DC}\ph{\second{.35}} & \cellcolor[HTML]{B8E1D4}\ph{\second{.46}} & \cellcolor[HTML]{D8EDE2}\ph{\second{.28}} & \ph{\second{.30}} \\
\famchip{foundation}\; RETFound & MAE encoder + UPerNet & ViT-L & \cellcolor[HTML]{D8EDE2}\ph{\best{.28}} & \cellcolor[HTML]{99D5CC}\ph{\best{.59}} & \cellcolor[HTML]{68C1C2}\ph{\best{.77}} & \cellcolor[HTML]{97D4CB}\ph{\best{.60}} & \ph{\best{.56}} & \cellcolor[HTML]{EBF4EB}\ph{\best{.12}} & \cellcolor[HTML]{C6E6DA}\ph{\best{.38}} & \cellcolor[HTML]{B3DFD2}\ph{\best{.48}} & \cellcolor[HTML]{D4ECE0}\ph{\best{.30}} & \ph{\best{.32}} \\
\midrule
\famchip{fundus}\; M2MRF & reference, Table~\ref{tab:main} & -- & \ph{.45} & \ph{.63} & \ph{.87} & \ph{.68} & \ph{.66} & \ph{.15} & \ph{.41} & \ph{.53} & \ph{.39} & \ph{.37} \\
\famchip{auto}\; nnU-Net & reference, Table~\ref{tab:main} & -- & \ph{.50} & \ph{.67} & \ph{.87} & \ph{.62} & \ph{.66} & \ph{.16} & \ph{.50} & \ph{.61} & \ph{.45} & \ph{.43} \\
\bottomrule
\end{tabular}
\begin{tablenotes}[flushleft]\scriptsize
\item $^\ddagger$The zero-shot block is summarised by the any-lesion AUPR (union of the four classes against the pixelwise maximum of the maps); automatic masks carry no class, so their per-class cells are empty, and SAM~3's per-class cells are its own text prompts. The any-lesion AUPR of the references on the same test sets (IDRiD / DDR): M2MRF -- / --; nnU-Net 0.78 / 0.57 (-- where the run predates the metric).
\item $^\dagger$Each automatic mask is assigned the class it overlaps most in the ground truth, an upper bound for class-agnostic proposals. Scores of automatic masks use the predicted IoU.
\phnote
\end{tablenotes}
\end{threeparttable}
\end{adjustbox}
\end{table*}


% ============================================================== R6
\subsection{Effects of annotation granularity and training source}
The sources differ in how lesions were drawn before they differ in anything a model does (Fig.~\ref{fig:sources}a). In the fine-grained sources an image carries many small components; in the coarse and region-consensus sources it carries few and large ones, which for the same retina means that several lesions and the tissue between them were enclosed by one outline.

\begin{simtext}
To separate the effect of the training source from the effect of the model, the five representative models were also trained on the standard-field TJDR images and on FGADR, at the pixel budget of DDR and with one seed, which together with the main tables gives a matrix of four training sources by seven test sets (Fig.~\ref{fig:sources}b,d). Every test set was served best by its own source. Training elsewhere cost \rn{cross.mean} \mAUPR{} on average and up to \rn{cross.max} (Fig.~\ref{fig:sources}c); models trained on the mixed labels of FGADR lost \rn{cross.fgadr.fine} on the three fine-grained sources.

We then added mixed or coarse labels to fine-label training while keeping the number of image passes fixed, for U-Net and M2MRF with one seed; differences smaller than twice the median seed standard deviation of the main table are reported without interpretation (Fig.~\ref{fig:sources}e and Table~\ref{tab:factors}). Relative to DDR alone (\rn{e4.ref} on the IDRiD test set), adding FGADR changed \mAUPR{} by \rn{e4.fgadr}; adding only 383 FGADR images, as many as DDR has, by \rn{e4.fgadr383}, which separates the effect of more images from the effect of a different label style. Adding Retinal-Lesions changed it by \rn{e4.retles}, and training on Retinal-Lesions or FGADR alone by \rn{e4.retlesonly} and \rn{e4.fgadronly}.
\end{simtext}

\begin{figure}[t]
\centering
\widefig{fig6_sources}
\caption{\textbf{Effects of annotation granularity and training source.} \textbf{a}, Annotation style: lesions per image against the median lesion diameter relative to the field-of-view diameter, per lesion class and source. \textbf{b}, \mAUPR{} for each training source (rows) and test set (columns), mean of five representative models, seed 0; outlined cells are internal tests. \textbf{c}, Difference to training on the source of the test set. \textbf{d}, The matrix of \textbf{b} for each model. \textbf{e}, Training sets of different granularity at an equal number of image passes, tested on two fine-grained test sets (U-Net and M2MRF, seed 0).}
\label{fig:sources}
\end{figure}

% ============================================================== R7
\subsection{Effects of lesion size, input resolution and training amount}
Lesion size is the property that separates the four classes most clearly (Fig.~\ref{fig:scale}a). On the training images the median diameter of a microaneurysm is 0.5\% of the field-of-view diameter and that of a hemorrhage 1.1\%; we froze these two values, computed on training images only, as the boundaries between small, medium and large lesions before the main-table training. Resizing a whole image to a small input removes lesions from the training target itself: at 224 pixels \rn{vanish.IDRiD.224.MA}\% of the microaneurysms and \rn{vanish.IDRiD.224.EX}\% of the hard exudates of the IDRiD training set no longer cover a single pixel, and on DDR \rn{vanish.DDR.224.MA}\% and \rn{vanish.DDR.224.EX}\%; at 384 pixels the DDR figures are \rn{vanish.DDR.384.MA}\% and \rn{vanish.DDR.384.EX}\%, and on the 1,440-pixel canvas of the unified protocol essentially none are lost (Fig.~\ref{fig:scale}b).

\begin{simtext}
Lesion-level recall followed size for every model (Fig.~\ref{fig:scale}c). On DDR the median recall over models was \rn{recall.DDR.small.median} for small and \rn{recall.DDR.large.median} for large lesions, and the best model detected \rn{recall.DDR.small.best} of the small ones. Input resolution acted mainly through this channel. For M2MRF and for the RETFound encoder with a UPerNet head, \mAUPR{} on DDR rose from \rn{e7.aupr.512} at a long side of 512 pixels to \rn{e7.aupr.1440} at 1,440 and \rn{e7.aupr.native} with native-resolution patches (mean of the two models, one seed), while the recall of small lesions rose from \rn{e7.small.512} to \rn{e7.small.1440} (Fig.~\ref{fig:scale}d).

More training images helped up to the size of the available sources. With 10\%, 50\% and 100\% of the DDR training set, \mAUPR{} on the DDR test set was \rn{e8.ddr.10}, \rn{e8.ddr.50} and \rn{e8.ddr.100} (U-Net and M2MRF, one seed, equal number of image passes), and merging the three fine-grained sources IDRiD, DDR and standard-field TJDR gave \rn{e8.ddr.all}. A source added to the training set raised the score on its own test set (TJDR standard field: \rn{e8.tjdr.100} before and \rn{e8.tjdr.fine} after it was added), which is the gain of in-source data seen in the matrix above and not a gain in generalization (Fig.~\ref{fig:scale}e). Exchanging the loss of U-Net with the architecture fixed moved \mAUPR{} by at most \rn{e9.range}, except for the per-image Dice term of the registered recipe, which is the collapse described above (Extended Data Fig.~8).
\end{simtext}

\begin{figure}[p]
\centering
\widefig{fig7_scale}
\caption{\textbf{Effects of lesion size, input resolution and training amount.} \textbf{a}, Lesion diameter relative to the field-of-view diameter on the training images; dashed lines, the frozen boundaries between small, medium and large lesions. \textbf{b}, Share of connected components that vanish from the training target when the whole image is resized to the given size. \textbf{c}, Lesion-level recall by lesion size and lesion-level precision on the internal tests, at the threshold chosen on the validation set. \textbf{d}, \mAUPR{} and recall of small lesions on DDR against the long side of the canvas (M2MRF and RETFound with UPerNet, seed 0). \textbf{e}, \mAUPR{} against the fraction of the DDR training set and after the fine-grained sources are merged in (U-Net and M2MRF, seed 0); the number of image passes is the same in every setting.}
\label{fig:scale}
\end{figure}

\begin{table*}[t]
\centering\tablefont
\caption{\textbf{One factor at a time.} One or two representative models are retrained with a single factor changed and everything else fixed to the unified protocol (seed 0). Each cell is the value followed by the change against the reference row of its block; -- marks a model the block does not use. The last column averages the models of the block; its shading follows the change. Differences smaller than twice the median seed standard deviation of the main table are not interpreted.}
\label{tab:factors}
\begin{adjustbox}{max width=\textwidth}
\begin{threeparttable}
\begin{tabular}{@{}l ccc c@{}}
\toprule
\colh{Setting} & \famchip{general}\; U-Net & \famchip{fundus}\; M2MRF & \famchip{foundation}\; RETFound+UPer. & \colh{Mean} \\
\midrule
\grouprow{5}{Annotation granularity of the training set (E4) --- mAUPR on IDRiD test, equal image passes; U-Net and M2MRF}
\quad DDR \tier{fine} {\scriptsize\color{inkmuted}(ref.)} & \ph{.36} & \ph{.54} & -- & \ph{.45} \\
\quad DDR + FGADR, 383 images \tier{fine}\tier{mixed} & \ph{.36} \same{\ph{.00}} & \ph{.54} \up{\ph{.01}} & -- & \cellcolor[HTML]{EFEFEC}\ph{.45} {\scriptsize \ph{+.00}} \\
\quad DDR + FGADR \tier{fine}\tier{mixed} & \ph{.38} \up{\ph{.02}} & \ph{.59} \up{\ph{.06}} & -- & \cellcolor[HTML]{E2E7EC}\ph{.48} {\scriptsize \ph{+.04}} \\
\quad DDR + Retinal-Lesions \tier{fine}\tier{coarse} & \ph{.32} \down{\ph{.04}} & \ph{.52} \down{\ph{.02}} & -- & \cellcolor[HTML]{EFE4E1}\ph{.42} {\scriptsize \ph{$$-$$.03}} \\
\quad FGADR \tier{mixed} & \ph{.27} \down{\ph{.09}} & \ph{.46} \down{\ph{.08}} & -- & \cellcolor[HTML]{EECFCA}\ph{.36} {\scriptsize \ph{$$-$$.09}} \\
\quad Retinal-Lesions \tier{coarse} & \ph{.21} \down{\ph{.15}} & \ph{.35} \down{\ph{.18}} & -- & \cellcolor[HTML]{EBB0A8}\ph{.28} {\scriptsize \ph{$$-$$.17}} \\
\grouprow{5}{Input resolution, long side in pixels (E7) --- mAUPR on DDR test; M2MRF and RETFound + UPerNet}
\quad 512 & -- & \ph{.23} \down{\ph{.19}} & \ph{.23} \down{\ph{.12}} & \cellcolor[HTML]{EBB3AC}\ph{.23} {\scriptsize \ph{$$-$$.16}} \\
\quad 1,440 {\scriptsize\color{inkmuted}(ref.)} & -- & \ph{.42} & \ph{.35} & \ph{.39} \\
\quad native-resolution patches & -- & \ph{.41} \down{\ph{.01}} & \ph{.34} \down{\ph{.01}} & \cellcolor[HTML]{F0EBE8}\ph{.38} {\scriptsize \ph{$$-$$.01}} \\
\grouprow{5}{Input resolution (E7) --- lesion-level recall of small lesions on DDR test}
\quad 512 & -- & \ph{.05} \down{\ph{.29}} & \ph{.11} \down{\ph{.20}} & \cellcolor[HTML]{E99289}\ph{.08} {\scriptsize \ph{$$-$$.24}} \\
\quad 1,440 {\scriptsize\color{inkmuted}(ref.)} & -- & \ph{.34} & \ph{.31} & \ph{.32} \\
\quad native-resolution patches & -- & \ph{.39} \up{\ph{.05}} & \ph{.29} \down{\ph{.01}} & \cellcolor[HTML]{E9EBEC}\ph{.34} {\scriptsize \ph{+.02}} \\
\grouprow{5}{Training amount (E8) --- mAUPR on DDR test; U-Net and M2MRF}
\quad DDR, 10\% (38 images) & \ph{.18} \down{\ph{.16}} & \ph{.23} \down{\ph{.17}} & -- & \cellcolor[HTML]{EBB1AA}\ph{.20} {\scriptsize \ph{$$-$$.16}} \\
\quad DDR, 50\% (192) & \ph{.29} \down{\ph{.05}} & \ph{.35} \down{\ph{.04}} & -- & \cellcolor[HTML]{EFDEDA}\ph{.32} {\scriptsize \ph{$$-$$.05}} \\
\quad DDR, 100\% (383) {\scriptsize\color{inkmuted}(ref.)} & \ph{.33} & \ph{.39} & -- & \ph{.36} \\
\quad + IDRiD + TJDR standard field (599) & \ph{.37} \up{\ph{.04}} & \ph{.44} \up{\ph{.04}} & -- & \cellcolor[HTML]{E0E7EC}\ph{.40} {\scriptsize \ph{+.04}} \\
\grouprow{5}{Loss, architecture fixed (E9) --- mAUPR on DDR test; U-Net}
\quad cross-entropy {\scriptsize\color{inkmuted}(ref.)} & \ph{.30} & -- & -- & \ph{.30} \\
\quad batch Dice & \ph{.34} \up{\ph{.04}} & -- & -- & \cellcolor[HTML]{E0E7EC}\ph{.34} {\scriptsize \ph{+.04}} \\
\quad cross-entropy + batch Dice (unified protocol) & \ph{.34} \up{\ph{.04}} & -- & -- & \cellcolor[HTML]{DFE6EC}\ph{.34} {\scriptsize \ph{+.04}} \\
\quad cross-entropy + per-image Dice (registered recipe) & \ph{.12} \down{\ph{.18}} & -- & -- & \cellcolor[HTML]{EBABA3}\ph{.12} {\scriptsize \ph{$$-$$.18}} \\
\bottomrule
\end{tabular}
\begin{tablenotes}[flushleft]\scriptsize
\item Granularity of the labels used for training: \tier{fine} fine \quad \tier{mixed} mixed \quad \tier{coarse} coarse. In the training-amount block the number of image passes is the same in every row. The FGADR row of the granularity block is the FGADR-trained model of the cross-source experiment.
\phnote
\end{tablenotes}
\end{threeparttable}
\end{adjustbox}
\end{table*}


% ============================================================== R8
\subsection{Computational cost}
\begin{simtext}
All models were trained on single 40\,GB cards, one task per card. Because the training amount is fixed in pixels, the cost of a model under the unified protocol reflects its throughput and its patch size: one seed on IDRiD took from \rn{cost.min} GPU hours (\rn{cost.min.model}) to \rn{cost.max} (\rn{cost.max.model}), and a DDR seed about four times as long, the slowest being the models that process small patches with sliding windows (FCT, 72 hours on DDR). Accuracy was not related to cost or to parameter count (Extended Data Fig.~10 and Supplementary Table~7).
\end{simtext}

\FloatBarrier
